\chapter{Background and Related Work}
\label{ch:background}

\section{DDoS taxonomy in data-center networks}
\label{sec:ddos-taxonomy}

DDoS attacks relevant to the DCN setting fall into three broad families.
\emph{Volumetric} floods (UDP, ICMP, amplification) aim to saturate links
at the edge or aggregation layer. \emph{Protocol} exhaustion attacks
(SYN flood, TCP state exhaustion) target control-plane state of middleboxes
and servers. \emph{Application-layer} attacks (HTTP flood, Slow-rate)
mimic legitimate traffic at high rate or low rate and are notoriously hard
to separate from flash-crowd events. CICDDoS2019 remains the most widely
cited public dataset covering the first two families; application-layer
samples are present as labels but with incomplete flow data, which
constrains empirical studies and motivates our multi-class design (see
\S\ref{sec:dataset-limits}).

\section{NFV, SFC and OpenFlow}

Network Function Virtualization (NFV), standardised by ETSI, decouples
packet-processing functions from specialised hardware. Service Function
Chaining (SFC) composes such VNFs into ordered chains; OpenFlow (and its
OVS implementation) provides the south-bound mechanism to steer traffic
through a chain. For DDoS mitigation the relevant VNFs are rate-limiters,
scrubbers, flow analysers, and black-holing functions, each with distinct
boot latency and resource footprint.

\section{ONAP MANO closed loop}
\label{sec:onap}

ONAP is the Linux Foundation reference orchestrator. The relevant
sub-components for closed-loop DDoS response are:
\begin{itemize}
  \item \textbf{SO} (Service Orchestrator) --- instantiates/terminates VNFs
        through the \texttt{/onap/so/infra/serviceInstantiation/v7/} REST API.
  \item \textbf{CLAMP} (Closed Loop Automation Management Platform) --- owns
        closed-loop templates and pushes policies.
  \item \textbf{Policy Framework} --- evaluates Drools rules, issues
        operational policies.
  \item \textbf{DMaaP} --- the Kafka-based message bus that carries AI
        decisions to CLAMP/Policy.
\end{itemize}
\padonap{} emits a canonical \texttt{AIOutputPayload} JSON on DMaaP; the
remaining closed-loop path is ONAP-native.

\section{ML models for network security}

Three model families dominate the literature. \emph{Tree ensembles}
(XGBoost~\citep{chen2016xgboost}, LightGBM, Random Forests) offer strong
accuracy on tabular flow features and sub-millisecond inference. \emph{Deep
sequence models} (LSTM, Transformer~\citep{vaswani2017attention}) capture
temporal structure and enable forecasting. \emph{Graph neural networks}
model host-to-host interaction but require flow-level graph construction.
\padonap{} combines XGBoost (detection) with Transformer+LSTM (forecasting)
to exploit both families.

\section{Forecasting horizons and the proactive gap}

A forecaster that predicts attack probability at time \(t{+}h\) permits
\emph{pre-positioning} --- instantiating VNFs \emph{before} the attack
peaks. The predictable lead-time is bounded above by the forecast horizon
\(h\) and below by the VNF boot latency. Choosing \(h\) too short collapses
to reactive behaviour; choosing \(h\) too long degrades AUC beyond useful.
We pick four horizons \(h \in \{30,60,90,120\}\) s and show (§\ref{sec:forecast})
that AUC degrades monotonically from \(0.988\) to \(0.971\).

\section{Fat-tree DCNs}
We evaluate on a Mininet fat-tree with \(k{=}4\) (20 switches, 16 hosts),
following the classical Al-Fares design, which gives three pod levels
(edge, aggregation, core). This allows us to measure the cost of
installing SFC rules on paths that traverse 1, 2, or 3 pod levels ---
a measurement absent from reactive baselines that assume single-switch
deployment.

\section{Related work comparison}
\label{sec:related}

Table~\ref{tab:related} compares eight recent systems across six axes
relevant to this thesis. The only system combining \emph{all four pillars}
(ML detection, multi-horizon forecast, ONAP-compliant orchestration, and
LP-based multi-tenant SLA) is \padonap{}.

\begin{table}[H]
\centering
\footnotesize
\caption{Related-work comparison matrix. $\bullet$ = supported, $\circ$ = partial, blank = absent.}
\label{tab:related}
\begin{tabular}{@{}lcccccc@{}}
\toprule
System & ML det. & Forecast & MANO & Multi-tenant SLA & Proactive & DCN eval. \\
\midrule
TST-DDoS (IEEE Access 2025)        & $\bullet$ &            &            &            &            &            \\
NetProbe (ICDCN 2025)               & $\bullet$ &            & $\circ$    &            & $\circ$    & $\bullet$  \\
FlowGuard (NDSS 2022)               & $\bullet$ & $\circ$    &            &            & $\circ$    &            \\
Lightweight ML (Sci.~Rep.~2025)     & $\bullet$ &            &            &            &            &            \\
DAWN (Springer LNCS 2025)           & $\bullet$ &            & $\circ$    &            &            & $\circ$    \\
Transformer+TCN (PLOS One 2025)     & $\bullet$ & $\circ$    &            &            & $\circ$    &            \\
Latency-aware NFV-6G (MDPI 2024)    & $\circ$   &            & $\bullet$  & $\circ$    &            & $\circ$    \\
Cross-layer SDN/NFV (MDPI 2025)     & $\bullet$ &            & $\bullet$  &            &            & $\bullet$  \\
\textbf{\padonap{} (ours)}          & \textbf{$\bullet$} & \textbf{$\bullet$} & \textbf{$\bullet$} & \textbf{$\bullet$} & \textbf{$\bullet$} & \textbf{$\bullet$} \\
\bottomrule
\end{tabular}
\end{table}

\section{Research gap}
From Table~\ref{tab:related}, three gaps motivate this work:
(G1) no prior system pairs multi-horizon forecasting with MANO-compliant
lifecycle management; (G2) no prior system ties mitigation tier to an
LP-backed tenant SLA guarantee; (G3) graceful handling of OOD traffic
(application-layer and ICMP variants absent from training) is rarely
measured and, when measured, often yields over-escalation.
